\section{Preliminaries}\label{sec:prelims}
%\kristof{$-$ and $\setminus$ are both used throughout the paper. We should unify the notation.}

%\knote{Consistently use ``for every''. Avoid ``for all'', ``for each'', etc. I have already done this in the intro and results section. }

We denote the sets of 
%\emph{reals}, \emph{non-negative reals}, \emph{positive reals}, 
\emph{integers}, \emph{non-negative integers}, and \emph{positive integers} by 
%$\mathbb{R}$, $\mathbb{R}_{\geq 0}$, $\mathbb{R}_+$, 
$\mathbb{Z}$, $\mathbb{Z}_{\geq 0}$, and $\mathbb{Z}_+$ respectively. For a positive integer $n$, we use $[n]$ to denote the set $\{1, 2, \ldots, n\}$. %We denote the value of the \emph{indicator function} for a statement $S$ by $\indicator_S$, i.e. $\indicator_S=1$ of $S$ is true and $0$ otherwise. \knote{State only what is used.}
Let $V$ be a finite set. We use $2^{V}$ to denote all subsets of $V$. 
%The \emph{characteristic vector} of a set $X\subseteq V$ is denoted by $\chi_X \in \{0, 1\}^V$ where $\chi_X(u)=1$ if $u\in X$ and $0$ otherwise for every $u\in V$. 
Let $X, Y\subseteq V$. We will denote $X\setminus Y$ by $X-Y$ and $X\cup Y$ by $X+Y$. If $Y$ consists of a single element $y$, then $X-\{y\}$ and $X+\{y\}$ are abbreviated as $X-y$ and $X+y$, respectively. 
Let $\calF\subseteq 2^V$ be a family of subsets of $V$. 
%A set $T\subseteq V$ is a \emph{transversal} for the family $\calF$ if each set in $\calF$ has non-empty intersection with the set $T$. 
The family $\calF$ is 
a \emph{disjoint family} if $X\cap Y=\emptyset$ for every $X, Y\in \calF$, 
a \emph{chain family} if either $X\subseteq Y$ or $Y\subseteq X$ for every $X, Y\in \calF$, and a \emph{laminar family} if either $X\cap Y=\emptyset$ or $X\subseteq Y$ or $Y\subseteq X$ for every $X, Y\in \calF$. 
We recall that the size of a laminar family $\mathcal{F}\subseteq 2^V$ is at most $2|V|$. 
For a subset $U\subseteq V$, we define the projection of $\mathcal{F}$ to $U$ as   $\mathcal{F}|_U:=\{X\cap U: X\in \mathcal{F}\}-\{\emptyset\}$ (we will use the convention that $\mathcal{F}|_U$ is a set and not a multi-set). 
%We will be interested in the \emph{projection} of a laminar family
%, i.e., the restriction of a laminar family $\calL \subseteq 2^V$ to a ground set $V - \calZ$ for some $\calZ\subseteq V$. 
The next lemma shows that the projection of a laminar family is also laminar and the size of the projection is comparable to that of the original family. We give a proof in \Cref{appendix:sec:uncrossing-properties:projection-of-laminar-families}.  %\knote{State only what is used.}\shubhang{I think we use everything except the set addition notation --- whether this is used anywhere still has to be verified}
%$\calL$.

\begin{restatable}{lemma}{projectionLaminarFamily}\label{lem:CoveringAlgorithm:projection-laminar-family}
    Let $\calL \subseteq 2^V$ be a laminar family and $\calZ \subseteq V$ be a subset of elements. 
    Let $\calL':=\calL|_{V-\zeros}$. 
    %Let $\calL' \coloneqq  \{X - \calZ : X \in \calL\} - \{\emptyset\}$. 
    Then, the family $\calL'$ is a laminar family and $|\calL| \leq |\calL'| + 3|\calZ|$.
\end{restatable}

%We prove \Cref{lem:CoveringAlgorithm:projection-laminar-family} in \Cref{appendix:sec:uncrossing-properties:projection-of-laminar-families}.


%\paragraph{Set Functions.} 
Let $p: 2^{V} \rightarrow \mathbb{R}$ be a set function over a finite ground set $V$. 
%We will denote the maximum function value of $p$ by $K_p$, i.e., $K_p\coloneqq \max\{p(X): X\subseteq V\}$ \knote{Is $K_p$ used at all?}.
The function $p$ is \emph{monotone} if  $p(X)\le p(Y)$ for every $X\subseteq Y\subseteq V$. 
%If the function $p$ is locally supermodular 
If $p(X)+p(Y)\le p(X\cap Y)+p(X\cup Y)$
for every $X,Y\subseteq V$, then the function $p$ is said to be \emph{supermodular}. 
%If the reverse of the same inequality holds for every $X,Y\subseteq V$, then the function $p$ is said to be \emph{submodular}. 
If $-p$ is supermodular, then the function $p$ is said to be \emph{submodular}. 
We recall that for a hypergraph $(H=(V, E), w: E\rightarrow \Z_+)$, the cut function $d_{(H, w)}: 2^V\rightarrow \Z_{\ge 0}$ is symmetric submodular and the coverage function $b_{(H, w)}: 2^V\rightarrow \Z_{\ge 0}$ is monotone submodular \cite{Frank-book}. 
%A set $X\subseteq V$ is said to be a \emph{$p$-maximizer} if $p(X) = \max\{p(Z): Z\subseteq V\}$. We use the notation $Z \coloneqq  \arg\max\{p(Z): Z\subseteq V\}$ to denote an arbitrary set $Z\subseteq V$ such that $p(Z)\ge p(X)$ for every $X\subseteq V$ (similarly for minimizers of the function $p$).
For a subset $Z \subseteq V$, the \emph{contracted function} $\functionContract{p}{Z} : 2^{V - Z}\rightarrow\Z$ is defined as $\functionContract{p}{Z}(X) \coloneqq  \max\{p(X\cup R) : R \subseteq Z\}$. We note that if $p$ is a skew-supermodular function and $Z\subseteq V$, then $\functionContract{p}{Z}$ is skew-supermodular. Moreover, if $p$ is a symmetric function and $Z\subseteq V$, then $\functionContract{p}{Z}$ is symmetric. 
%For a function $m: V\rightarrow \R$ and a subset $Z\subseteq V$, we define the \emph{restricted function} $\functionRestrict{m}{Z} : V-Z \rightarrow \R$ as $\functionRestrict{m}{Z}(u) \coloneqq  m(u)$ for every $u\in V - Z$.  
%\knote{State only what is used.}

We recall that in our algorithmic problems, we have access to the input set function $p$ via \functionMaximizationOracleStrongCover{p}. 
%We describe an additional oracle that  will simplify our proofs. We will show that both oracles defined below can be implemented in strongly polynomial time using \functionMaximizationOracleStrongCover{p}. 
We describe an additional oracle that will simplify our proofs. We will show that the oracle defined below can be implemented in strongly polynomial time using \functionMaximizationOracleStrongCover{p}. 
\begin{definition}[$p$-maximization oracle]
    Let $p: 2^V\rightarrow \Z$ be a set function. 
    \functionMaximizationOracle{p}$((G_0, c_0), S_0, T_0)$ takes as input a hypergraph $(G_0=(V, E_0), c_0: E_0\rightarrow \mathbb{Z}_{+})$  
    and disjoint sets $S_0, T_0\subseteq V$, and returns a tuple $(Z, p(Z))$, where $Z$ is an optimum solution to the following problem:
    \begin{align*}
        \max & \left\{p(Z)-b_{(G_0, c_0)}(Z): S_0\subseteq Z\subseteq V-T_0\right\}. \tag{\functionMaximizationOracle{p}}\label{tag:oracle-main}
        \end{align*} 
\end{definition}
%\knote{If we want strongly polynomial time algorithms, then we can query the above oracles only on hypergraphs $(G_0, c_0)$, where $|E_0|=$ poly$(|V|)$ and $c_0$ is also bounded by ...}
The evaluation oracle for a function $p$ can be implemented using one query to  
%\functionMaximizationEmptyOracle{p};  
%\functionMaximizationEmptyOracle{p} can be implemented using one query to 
\functionMaximizationOracle{p} where the input hypergraph $(G_0, c_0)$ is the empty hypergraph and $S_0 := Z$, $T_0 := V - Z$. The next lemma shows that \functionMaximizationOracle{p} can be implemented using at most $|V|+1$ queries to \functionMaximizationOracleStrongCover{p} where the hypergraphs used as input to  \functionMaximizationOracleStrongCover{p} have size of the order of the size of the hypergraph input to \functionMaximizationOracle{p}. We give a proof in \Cref{sec:wc-oracle-from-sc-oracle}.  

\begin{restatable}{lemma}{lemmaWCoraclefromSCoracle}\label{lem:Preliminaries:wc-oracle-from-sc-oracle}
Let $p: 2^V\rightarrow \Z$ be a set function, $(G = (V, E), w:E\rightarrow\Z_+)$ be a hypergraph, and $S, T \subseteq V$ be disjoint sets. Then, $\functionMaximizationOracle{p}((G, w), S, T)$ can be implemented to run in $O(|V|(|V|+|E|))$ time using at most $|V|+1$ queries to $\functionMaximizationOracleStrongCover{p}$.
The run-time includes the time to construct the hypergraphs used as input to the queries to \functionMaximizationOracleStrongCover{p}. Moreover, each query to \functionMaximizationOracleStrongCover{p} is on an input hypergraph $(G_0, c_0)$ that has at most $|V|$ vertices and $|E|$ hyperedges.
\end{restatable}


For two hypergraphs $(G=(V, E_G), c_G:E_G\rightarrow \Z_+)$ and $(H=(V, E_H), w_H: E_H\rightarrow \Z_+)$ on the same vertex set $V$, we recall that the hypergraph $(G+H=(V, E_{G+H}), c_G+w_H)$ is the hypergraph with vertex set $V$ and hyperedge set $E_{G+H}\coloneqq E_G\cup E_H$ with the weight of every hyperedge $e\in E_G\cap E_H$ being $c_G(e) + w_H(e)$, the weight of every hyperedge $e\in E_G\setminus E_H$ being $c_G(e)$,  and the weight of every hyperedge $e\in E_H\setminus E_G$ being $w_H(e)$. 